\section{Exactness and residual geometry}\label{sec:geometry}
Let $X$ be a nonempty compact subset of a finite-dimensional Euclidean
space. Let $f:X\to\R$ and $r:X\to\R^m$ be continuous, and suppose
$F=\{x\in X:r(x)=0\}$ is nonempty. Fix a norm on $\R^m$ and write
\begin{equation}\label{eq:model}
 v=\min_{x\in F}f(x),\qquad
 L_\rho(\lambda)=\min_{x\in X}\{f(x)+\lambda^Tr(x)+\rho\norm{r(x)}\},
 \qquad D_\rho=\sup_{\lambda\in\R^m}L_\rho(\lambda),
\end{equation}
where $\rho\ge0$. The set $X$ retains the \emph{native constraints};
only $r(x)=0$ is relaxed. Weak duality gives $L_\rho(\lambda)\le D_\rho\le v$.
We call $\rho$ \emph{dual-value exact} if $D_\rho=v$, and exact at a
specified multiplier $\lambda$ if $L_\rho(\lambda)=v$. The latter includes
attainment of the dual supremum. We call $(\rho,\lambda)$
\emph{minimizer-set exact} if the minimizers in $L_\rho(\lambda)$ are
exactly $\argmin_F f$. The distinctions are necessary even on compact sets.

Set $\rho_*=\inf\{\rho\ge0:D_\rho=v\}$, with $\inf\varnothing=+\infty$.
For a fixed multiplier, the corresponding value threshold is
\begin{equation}\label{eq:fixed}
 \rho_*(\lambda)=\max\left\{0,\sup_{x\in X\setminus F}
 \frac{v-f(x)-\lambda^Tr(x)}{\norm{r(x)}}\right\}.
\end{equation}
An empty supremum contributes no positive lower bound. This formula follows
by requiring every penalized objective value to be at least $v$.
A finite threshold is exact at that multiplier. Every strictly larger
coefficient is minimizer-set exact: its additional penalty is positive at
each infeasible point. In contrast, finite $\rho_*$ need not give an
attaining multiplier.

Primal convexification descriptions of augmented Lagrangian duals are
established \citep{BolandEberhard2015,FeizollahiAhmedSun2017}.
In particular, \citet[Theorem~1 of the inspected manuscript]{FeizollahiAhmedSun2017}
use a lifted convex hull under a relative-interior feasibility condition.
The following residual-space analogue for compact $X$ needs no
relative-interior condition. Its proof does not assume multiplier
attainment; we make no novelty claim for the convexification argument.

\begin{proposition}[Balanced mixtures]\label{prop:mixture}
For every $\rho\ge0$,
\begin{equation}\label{eq:mixture}
 D_\rho=\min\left\{\sum_{i=1}^p\theta_i
 [f(x_i)+\rho\norm{r(x_i)}]:
 \begin{array}{l}
 1\le p\le m+1,\ x_i\in X,\ \theta_i\ge0,\\
 \sum_i\theta_i=1,\quad\sum_i\theta_i r(x_i)=0
 \end{array}\right\}.
\end{equation}
\end{proposition}
\begin{proof}
The convex hull $K$ of
$\{(r(x),f(x)+\rho\norm{r(x)}):x\in X\}$ is compact.
Let $w=\min\{t:(0,t)\in K\}$. Averaging over a balanced mixture cancels
$\lambda$, so $D_\rho\le w$. For $a<w$, strict separation of $(0,a)$
from $K$ gives $u^Tz+\beta t>\beta a$ for every $(z,t)\in K$.
At $(0,w)$ this implies $\beta>0$. Compactness then gives
$L_\rho(u/\beta)>a$. Letting $a\uparrow w$ proves $D_\rho=w$, without
assuming multiplier attainment.

Carath\'eodory's theorem gives an optimal mixture with at most $m+2$
points. If it has more than $m+1$ positive weights, the columns
$(1,r(x_i))$ are linearly dependent. Both signs of a sufficiently small
weight perturbation along a dependence preserve feasibility.
Optimality forces its objective derivative to be zero. Increase the
perturbation until a weight vanishes; the value is unchanged. This
reduces the support to at most $m+1$.
\end{proof}
The average in \eqref{eq:mixture} contains the norms of the individual
residuals, not the norm of their average, which is zero.
The support bound is sharp: take one point with $(r,f)=(0,0)$ and
$m+1$ points with cost $-1$ and residuals
$e_1,\ldots,e_m,-\sum_j e_j$. At $\rho=0$, an optimal balanced mixture
must give weight $1/(m+1)$ to each of the latter points.

\begin{corollary}\label{cor:ratio}
Let $\mathcal B$ denote the balanced mixtures in \eqref{eq:mixture}, and
write $c_\mu=\sum_i\theta_i f(x_i)$ and
$s_\mu=\sum_i\theta_i\norm{r(x_i)}$. Then
\begin{equation}\label{eq:ratio}
 \rho_*=\max\left\{0,\sup_{\mu\in\mathcal B:s_\mu>0}
                  \frac{v-c_\mu}{s_\mu}\right\}.
\end{equation}
If this number is finite, it is dual-value exact.
\end{corollary}
\begin{proof}
A mixture with $s_\mu=0$ uses only feasible points of positive weight,
so $c_\mu\ge v$. The remaining inequalities in \eqref{eq:mixture} are
$\rho s_\mu\ge v-c_\mu$. They all hold at their finite supremum.
\end{proof}

\begin{proposition}[One relaxed equality]\label{prop:scalar}
Suppose $m=1$, use absolute value, and assume residuals of both signs occur.
Define
\[
 A_+=\sup_{r(x)>0}\frac{v-f(x)}{r(x)},\qquad
 A_-=\sup_{r(x)<0}\frac{v-f(x)}{-r(x)}.
\]
If both are finite, then
\begin{equation}\label{eq:scalar}
 \rho_*=\max\{0,(A_++A_-)/2\},\qquad
 L_\rho(\lambda)=v\ \Longleftrightarrow\
 A_+-\rho\le\lambda\le\rho-A_-.
\end{equation}
If either is infinite, no finite penalty is dual-value exact.
At multiplier zero the threshold is $\max\{0,A_+,A_-\}$.
\end{proposition}
\begin{proof}
The pointwise inequalities for positive and negative residuals give the
interval in \eqref{eq:scalar}. To see that an exact supremum cannot avoid
this interval, fix one point of each residual sign. Evaluating any sequence
$L_\rho(\lambda_j)\to v$ at these two points bounds $\lambda_j$ above
and below. Along a convergent subsequence, upper semicontinuity of
$L_\rho$ gives an attaining multiplier. Hence exactness requires and is
implied by a nonempty interval, giving all the claims.
\end{proof}
If residuals occur on only one side of zero, every balanced mixture is
supported on $F$, so $D_\rho=v$ for all $\rho\ge0$. Attainment can fail:
for $X=[0,1]$, $r(x)=x$, and $f(x)=-\sqrt{x}$, every finite multiplier
and penalty give $L_\rho(\lambda)<0=v$. Appendix~\ref{app:attainment}
records the same distinction in a mixed-integer quadratic example from
the literature.

\begin{lemma}[A bound from integer slices]\label{lem:slices}
Suppose $X$ is the union of finitely many nonempty compact slices $X_z$,
$f\ge L$ on $X$, and $v\le U$. On every slice meeting $F$, suppose
there is $\lambda_z$ such that
\[
 f(x)+\lambda_z^Tr(x)\ge v_z:=\min_{X_z\cap F}f,
 \qquad \norm{\lambda_z}_*\le M\quad(x\in X_z),
\]
where $\norm{\cdot}_*$ is the dual norm. Suppose every other slice has
$\min_{X_z}\norm{r}\ge\delta>0$. Then
\begin{equation}\label{eq:slicebound}
 \rho\ge\max\{M,(U-L)/\delta\}
\end{equation}
is exact at multiplier zero. Any strictly larger coefficient is
minimizer-set exact. If all slices meet $F$, omit the second term.
\end{lemma}
\begin{proof}
On a feasible slice, the dual-norm inequality gives
$f+\rho\norm r\ge f+\lambda_z^Tr\ge v_z\ge v$.
On an infeasible slice, $f+\rho\norm r\ge L+\rho\delta\ge U\ge v$.
A primal minimizer attains $v$. Strictly increasing $\rho$ excludes every
infeasible point from the minimizing set.
\end{proof}
Thus feasible-slice multiplier bounds and infeasible-slice separation
control different parts of the penalty. For convex slices, Slater regularity can supply
the former without giving a useful quantitative bound on the latter.
